% p.63 — الباب التاسع والعشرون في الاصطلاحات والتركيبات
\mspage{63}
% Layout: one boxed table of four columns forming two pairs (رقم + اصطلاح | مثال),
% read right pair first (1–21) then left pair (22–42).  The numbers are plain Urdu
% digits in the lithograph (not circled); transcribed as (n) before the term.
% Printed page number ۶۳ at top left.  Layout: the two pairs are set side by side (right
% pair first) in one `tables` row, as in the lithograph.
\bab{الباب التاسع والعشرون في الاصطلاحات والتركيبات}
\begin{tables}
\begin{sarf}{|r|r|}
(1) اسم & زَيْدٌ . دَارٌ\\
(2) فعل & نَامَ ضَرَبَ يَنْصُرُ\\
(3) حرف & فِيْ اِنْ لَمْ\\
(4) كلمة & زَيْدٌ ضَرَبَ فِيْ\\
(5) مفرد & زَيْدٌ ضَرَبَ فِيْ\\
(6) معرفة & زَيْدٌ\\
(7) نكرة & رَجُلٌ\\
(8) تثنية & زَيْدَانِ رَجُلَانِ\\
(9) جمع & زَيْدُوْنَ رِجَالٌ\\
(10) معرب & زَيْدٌ زَيْدًا زَيْدٍ\\
(11) مبني & هٰؤُلَاءِ\\
(12) مركب & فِي الدَّارِ\\
(13) مركب ناقص & أَحَدَ عَشَرَ\\
(14) مركب تام & زَيْدٌ فِي الدَّارِ\\
(15) جملة & نَامَ زَيْدٌ اِنْ نَامَ زَيْدٌ\\
(16) كلام & نَامَ زَيْدٌ\\
(17) حرف جر & فِي الدَّارِ\\
(18) حرف مشبه بالفعل & اِنَّ زَيْدًا قَائِمٌ\\
(19) \unclear{حرف مشبة بليس} & مَا زَيْدٌ قَائِمًا\\
(20) حرف استثناء & مَا جَاءَ الْقَوْمُ اِلَّا زَيْدًا\\
(21) حرف ندا & يَا زَيْدُ يَا عَبْدَ اللّٰهِ\\
\end{sarf}
\begin{sarf}{|r|r|}
(22) حرف ناصب & لَنْ يَذْهَبَ زَيْدٌ\\
(23) حرف جازم & لَمْ يَلِدْ وَلَمْ يُوْلَدْ\\
(24) حرف شرط & اِنْ تَذْهَبْ\\
(25) حرف جزاء & اِنْ تَذْهَبْ فَاذْهَبْ\\
(26) حرف نفي & مَا ضَرَبَ لَا يَضْرِبُ\\
(27) حرف تنبيه & أَلَا أُخْبِرُكَ بِخَيْرِ الْأَعْمَالِ\\
(28) لام تأكيد & لَيَضْرِبَنَّ زَيْدٌ\\
(29) حرف مصدرية & أُرِيْدُ أَنْ أُسْلِمَ\\
(30) حرف تحقيق & قَدْ جَاءَ نَصْرُ اللّٰهِ\\
(31) حرف تقليل & قَدْ يَجِيْئُ عَذَابُ اللّٰهِ\\
(32) حرف تحضيض & أَلَا تَنْزِلُ بِنَا\\
(33) حرف غائية & أَضْرِبُكَ حَتّٰى تَسْتَقِيْمَ\\
(34) حرف تنفيس & سَيَضْرِبُ زَيْدٌ . سَوْفَ يَضْرِبُ\\
(35) حرف استفهام & هَلْ هٰذَا زَيْدٌ\\
(36) حرف عطف & رَأَيْتُ زَيْدًا وَ خَالِدًا\\
(37) لام نهي & لَا تَضْرِبْ\\
(38) لام امر & لِيَضْرِبْ\\
(39) اسم اشارة & هٰذَا زَيْدٌ\\
(40) اسم موصول & هٰذَا غُلَامٌ الَّذِيْ فَرَّ أَمْسِ\\
(41) اسم ضمير & زَيْدٌ هُوَ عَالِمٌ\\
(42) فعل ناقص & كَانَ زَيْدٌ عَالِمًا\\
\end{sarf}
\end{tables}
